% !TEX root = relatorio.tex
%
% Capítulo 2 — Estado da Arte
%
\chapter{Estado da Arte} \label{cap:estadodaarte}

Este capítulo situa o Play4Change no panorama da aprendizagem adaptativa, da Inteligência Artificial na educação, da distinção entre aprendizagem formal e informal, do \textit{Retrieval-Augmented Generation}, da aprendizagem móvel e da autenticação \textit{passwordless}, terminando com uma síntese comparativa face às soluções existentes.

% -----------------------------------------------------------
\section{Aprendizagem Adaptativa e Sistemas de Tutoria Inteligente} \label{sec:its}
% -----------------------------------------------------------

A ideia de adaptar o ensino ao ritmo e ao perfil individual de cada aluno não é nova: em meados do século XX, Skinner propôs as primeiras ``máquinas de ensinar'', baseadas no reforço imediato da resposta correta~\cite{skinner1958}, evoluindo nas décadas seguintes para os \textit{Intelligent Tutoring Systems} (ITS) --- sistemas que modelam o estado de conhecimento do aluno e adaptam dinamicamente a instrução através de três módulos, domínio, estudante e pedagógico~\cite{anderson1985,vanlehn2011}, o segundo tipicamente implementado por \textit{Knowledge Tracing}~\cite{corbett1994,piech2015}.

No mercado atual, plataformas como Duolingo, Khan Academy e Coursera incorporam adaptatividade em graus distintos, mas partilham uma limitação comum: o conteúdo é estático e curado manualmente por equipas especializadas, não podendo ser adaptado por terceiros para domínios específicos --- é aqui que o Play4Change se distingue, permitindo que qualquer Administrador introduza novo conhecimento, automaticamente transformado em conteúdo pedagógico pela IA. (O estudo do próprio Duolingo que reporta resultados comparáveis a um semestre universitário~\cite{vesselinov2012} foi, refira-se, publicamente contestado por não ser revisto por pares nem ter grupo de controlo independente~\cite{krashen2019}.)

% -----------------------------------------------------------
\section{Inteligência Artificial na Educação} \label{sec:aied}
% -----------------------------------------------------------

A aplicação da Inteligência Artificial à educação (\textit{AI in Education}, AIEd) tem uma história de várias décadas, mas os LLM introduziram uma mudança qualitativa nas possibilidades desta área.

\subsection{O Problema dos 2 Sigma e a Escalabilidade da Tutoria}

Em 1984, Bloom demonstrou que a tutoria individual supera em dois desvios padrão (\textit{2 sigma}) o ensino tradicional em sala de aula --- um ganho pedagógico economicamente inviável a essa escala~\cite{bloom1984}. Os LLM reabriram este debate~\cite{kasneci2023}: um modelo disponível 24 horas, capaz de adaptar a explicação ao nível do utilizador e dar \textit{feedback} imediato e contextualizado, aproxima pela primeira vez esse ideal de uma escala de implementação massiva~\cite{bommasani2021}.

\subsection{Feedback Formativo e Aprendizagem Visível}

A investigação converge num ponto, independentemente da tecnologia: o \textit{feedback} imediato e formativo acelera a aprendizagem --- Hattie (2009) coloca-o entre os dez maiores aceleradores, com um tamanho de efeito $d = 0{,}73$~\cite{hattie2009}, ainda que a robustez estatística desta meta-análise tenha sido questionada~\cite{simpson2017}. O Play4Change incorpora este princípio ao fornecer \textit{feedback} explicativo imediatamente após cada desafio diário.

\subsection{Teoria da Carga Cognitiva}

A \textit{Cognitive Load Theory} de Sweller~\cite{sweller1994} distingue carga cognitiva \textit{intrínseca}, \textit{extrínseca} e \textit{germane}; o desafio curto, diário e por tópico do Play4Change é consistente com estes princípios, embora a inscrição simultânea em vários tópicos ativos (Secção~\ref{sec:atribuicao} do Capítulo~\ref{cap:desafios}) dilua parcialmente esse argumento para esse perfil de utilização --- uma tensão de desenho que a versão atual não resolve explicitamente.

% -----------------------------------------------------------
\section{Aprendizagem Formal, Informal e Móvel} \label{sec:formal-informal}
% -----------------------------------------------------------

A distinção entre aprendizagem formal e informal é central para compreender o posicionamento do Play4Change. A \textbf{aprendizagem formal} caracteriza-se por ser institucionalizada, estruturada e avalizada por certificações reconhecidas, com motivação frequentemente extrínseca~\cite{colley2003}. A \textbf{aprendizagem informal}, por contraste, é não estruturada, autodirigida e motivada pela necessidade imediata ou curiosidade intrínseca~\cite{cross2007}: Knowles demonstrou que os adultos aprendem melhor quando a aprendizagem resolve problemas reais e é controlada pelo próprio~\cite{knowles1984}, e Cross popularizou a estimativa --- originada em investigação anterior do Institute for Research on Learning --- de que 80\% do conhecimento prático usado no trabalho é adquirido informalmente~\cite{cross2007}. O Play4Change opera deliberadamente na interseção destes dois paradigmas: ancorado no referencial estruturado do DigComp~2.2, mas com um modelo de utilização --- desafio diário, sem horários fixos nem avaliações sumativas --- intrinsecamente informal e autodirigido.

O canal escolhido para essa interseção é o dispositivo móvel: a aprendizagem móvel (\textit{m-learning}) tira partido da portabilidade e ubiquidade destes dispositivos~\cite{traxler2007}, transformando a aprendizagem num processo contínuo integrado na vida quotidiana~\cite{laurillard2007} --- central para o Play4Change, cujo desafio diário se integra nos momentos em que o utilizador já usa o dispositivo. Com 8,9 mil milhões de subscrições móveis ativas no mundo~\cite{itu2023}, é também o canal mais inclusivo disponível, relevante para o objetivo do DigComp~2.2 de servir \textit{todos} os cidadãos~\cite{digcomp22}.

% -----------------------------------------------------------
\section{Retrieval-Augmented Generation} \label{sec:rag-arte}
% -----------------------------------------------------------

Os LLM apresentam uma limitação estrutural relevante para aplicações de domínio específico: o seu conhecimento é estático, limitado à data de treino, e caro de atualizar sem re-treino do modelo.
O \textit{Retrieval-Augmented Generation} (RAG) surgiu como resposta a esta limitação, combinando a capacidade generativa dos LLM com a consulta dinâmica de bases de conhecimento externas~\cite{lewis2020}.

\subsection{Arquitetura RAG}

Na sua formulação original, proposta por Lewis et al. (2020), o RAG funciona em duas fases~\cite{lewis2020}:

\begin{enumerate}
    \item \textbf{Fase de recuperação (\textit{retrieval}):} dado um \textit{query}, o sistema recupera os documentos ou fragmentos de texto (\textit{chunks}) mais relevantes de uma base de conhecimento indexada vetorialmente, tipicamente por similaridade cosseno entre \textit{embeddings}.
    \item \textbf{Fase de geração:} os \textit{chunks} recuperados são injetados no contexto do LLM, que gera uma resposta fundamentada nessa informação específica.
\end{enumerate}

Esta arquitetura reduz as \textit{hallucinations} (respostas factualmente incorretas geradas com confiança), permite que o sistema responda com base em conhecimento privado ou recente, e torna as respostas rastreáveis às fontes originais~\cite{gao2024}.

\subsection{RAG em Contextos Educativos}

A aplicação do RAG a sistemas educativos é uma área de investigação em crescimento acelerado: ancorar a geração de questões e \textit{feedback} em documentos fornecidos pelo formador --- em vez de no conhecimento genérico do modelo --- é particularmente valioso em contextos de formação especializada, onde a precisão e a especificidade do conteúdo são críticas~\cite{kasneci2023}.
No Play4Change, o RAG garante que os desafios gerados para cada Learner estão diretamente ancorados no conteúdo submetido pelo Administrador, e não em conhecimento genérico do modelo.

% -----------------------------------------------------------
\section{Autenticação \textit{Passwordless}} \label{sec:auth-arte}
% -----------------------------------------------------------

A autenticação por \textit{password} é há décadas o mecanismo dominante de controlo de acesso, mas a investigação em segurança é consistente: as \textit{passwords} são reutilizadas entre serviços, vulneráveis a força bruta e \textit{phishing}, e frequentemente esquecidas~\cite{bonneau2012}. Bonneau et al. (2012) avaliaram 35 esquemas alternativos segundo 25 critérios e concluem que não há vencedor universal --- os esquemas mais resistentes a ataques (tokens físicos, biometria) tendem a sacrificar os critérios de usabilidade em que as \textit{passwords} se destacam, pelo que a escolha correta depende de uma ponderação explícita ao contexto de utilização, não de uma superioridade abstrata de um esquema sobre outro~\cite{bonneau2012}.

Aplicando esse critério de ponderação ao Play4Change --- um desafio breve e diário, dirigido também a um perfil de menor afinidade tecnológica (Secção~\ref{sec:usab-qualitativa}) --- o \textit{magic link}~\cite{fido2019} destaca-se por eliminar por completo a fricção de memorização e por resolver a recuperação de acesso por construção (pedir um novo \textit{link}, sem perguntas de segurança), a um custo aceite e explícito: a segurança da conta passa a depender da segurança do email do Learner. Um esquema FIDO2/\textit{Passkeys} eliminaria essa dependência, mas exige suporte de \textit{hardware} e uma curva de adoção que não se justifica para este perfil de utilizador. As propriedades criptográficas concretas do fluxo são detalhadas na Secção~\ref{sec:auth-model}.

% -----------------------------------------------------------
\section{Síntese Comparativa} \label{sec:sintese}
% -----------------------------------------------------------

A Tabela~\ref{tab:comparativo} apresenta uma análise comparativa entre o Play4Change e as principais plataformas de referência identificadas no estado da arte, segundo os critérios mais relevantes para o posicionamento do projeto.

\begin{table}[h!]
\centering
\caption{Comparação entre o Play4Change e plataformas de referência}
\label{tab:comparativo}
\resizebox{\textwidth}{!}{%
\begin{tabular}{lccccc}
\hline
 & & & & \multicolumn{2}{c}{\textbf{Play4Change}} \\
\cline{5-6}
\textbf{Critério} & \textbf{Duolingo} & \textbf{Khan Academy} & \textbf{Coursera} & \textbf{Atual} & \textbf{Futuro} \\
\hline
Gamificação                            & \cmark & \pmark & \xmark & \pmark & \cmark \\
Adaptação à dificuldade do utilizador  & \cmark & \cmark & \xmark & \cmark & \cmark \\
Catálogo de conteúdo pré-definido      & \cmark & \cmark & \cmark & \xmark & \pmark \\
Reconhecimento com credenciais formais & \pmark & \xmark & \cmark & \xmark & \cmark \\
Comunidade e interação social          & \cmark & \pmark & \cmark & \xmark & \cmark \\
\hline
Modelo de desafio diário no móvel      & \cmark & \xmark & \xmark & \cmark & \cmark \\
Cliente móvel nativo                   & \cmark & \cmark & \cmark & \cmark & \cmark \\
Autenticação \textit{passwordless}     & \xmark & \xmark & \xmark & \cmark & \cmark \\
\hline
\multicolumn{6}{l}{\footnotesize \cmark\ Suportado \quad \xmark\ Não suportado \quad \pmark\ Parcialmente suportado}
\end{tabular}%
}
\end{table}

A análise revela um posicionamento diferenciado, mas expõe também limitações claras: nos critérios onde as plataformas estabelecidas lideram --- gamificação, reconhecimento formal por certificados, catálogo pré-definido e comunidade/interação social --- o Play4Change fica em desvantagem, por depender do conteúdo submetido pelo Administrador e não cobrir ainda estas dimensões sociais, reconhecidas como prioridades de trabalho futuro. Diferencia-se, ainda assim, pelo seu modelo de entrega --- desafio diário por tópico no canal móvel, com adaptação ao desempenho e autenticação \textit{passwordless} --- combinação que nenhuma das plataformas analisadas oferece (com a ressalva sobre múltiplas inscrições simultâneas já assinalada na Secção~\ref{sec:aied}).
